\chapter{Problem Statement}\label{ch:problem}

This chapter translates the motivations from Chapter~\ref{ch:introduction} and the literature survey from Chapter~\ref{ch:background} into a precise research problem.
While automated performance tuning has achieved significant maturity in the domain of specific applications, most notably Database Management Systems (DBMS)~\cite{qtune,ottertune} and big data frameworks, these approaches typically treat the underlying Operating System (OS) as a static black box. They focus on optimizing application-level parameters (e.g., buffer pool sizes, query caches) while neglecting the kernel configuration upon which the application relies.
This thesis addresses this gap by targeting the \emph{search-space} bottleneck in automated performance tuning of the Linux kernel itself. Unlike application-specific tuners, the approach focuses on the generic system call interfaces and kernel subsystems, making the solution applicable across diverse workloads.

\section{Problem Context}\label{sec:problem-context}

Modern services commonly run inside single-tenant virtual machines or containers, yet they inherit a general-purpose Linux kernel configuration.
Linux exposes thousands of runtime parameters (e.g., via \texttt{sysctl}) that control major subsystems such as networking, memory management, and I/O~\cite{kernel_parameters}.
These parameters are organized hierarchically, as illustrated in Figure \ref{fig:sysctl_tree}.

\begin{figure}[ht]
    \centering
    \tikzset{
      myfolder/.style={fill=blue!10, draw=blue!50, thick, rounded corners},
      myfile/.style={fill=white, draw=gray!50, rounded corners, font=\sffamily\small}
    }

    \begin{forest}
      for tree={
        font=\sffamily\bfseries,
        if n children=0{myfile}{myfolder},
        edge={draw=darkgray, thick}
      }
      [/proc/sys, 
        grow=270,               % 1. Make the root grow downwards
        parent anchor=south,
        l sep=8mm,              % Vertical drop before the columns start
        s sep=2mm,              % Horizontal space between the 4 main columns
        for children={
          for tree={
            folder,             % 2. Make all sub-levels use the vertical folder style
            grow'=0,            % Folder children grow to the right
            fit=band,
            s sep=4pt           % Vertical space between files
          },
          % 3. Override the edge path for the 4 main columns so they connect neatly
          child anchor=north,
          edge path={
            \noexpand\path[\forestoption{edge}]
              (!u.parent anchor) -- +(0,-4mm) -| (.child anchor)\forestoption{edge label};
          }
        }
        [fs
          [file-max]
          [pipe-max-size]
          [aio-max-nr]
        ]
        [kernel
          [pid\_max]
          [sched\_migration\_cost\_ns]
          [numa\_balancing]
        ]
        [net
          [core
            [busy\_poll]
            [netdev\_max\_backlog]
            [rmem\_max]
          ]
          [ipv4
            [tcp\_window\_scale]
            [tcp\_congestion\_control]
            [ip\_forward]
          ]
        ]
        [vm
          [swappiness]
          [dirty\_ratio]
          [min\_free\_kbytes]
        ]
      ]
    \end{forest}
    
    \caption{Hierarchical structure of the \texttt{/proc/sys} interface. This view highlights the critical subsystems for performance tuning: filesystem (\texttt{fs}), kernel scheduling (\texttt{kernel}), networking (\texttt{net}), and memory management (\texttt{vm}).}
    \label{fig:sysctl_tree}
\end{figure}

Tuning these parameters can yield significant performance benefits, but doing so reliably is difficult:
\begin{itemize}
    \item \textbf{Scale:} the configuration space is extremely large (thousands of knobs, many with wide ranges).
    \item \textbf{Opacity:} it is rarely clear which knobs influence a given application.
    \item \textbf{Interactions:} knobs interact with each other and with workload characteristics.
    \item \textbf{Dynamics:} ``best'' settings can change with input diversity, kernel version, and hardware.
\end{itemize}

Consequently, operators often either leave defaults unchanged or apply generic tuning recipes.
This leads to sub-optimal performance and increases operational risk.

\section{The Core Challenge: Search Space Explosion}\label{sec:problem-search-space}

Autotuning systems search for configurations that optimize an objective such as throughput, latency, or resource efficiency.
The literature shows that both Bayesian Optimization and Reinforcement Learning are common choices for such search problems~\cite{stun2022,autotune_survey}.
However, their effectiveness rapidly degrades as dimensionality grows.

Let $\mathcal{K}$ be the set of all available sysctl kernel knobs, and let $|\mathcal{K}| = d$.
A configuration is a vector $\mathbf{x} \in \mathcal{X}$ where
$\mathcal{X} = \mathcal{X}_1 \times \cdots \times \mathcal{X}_d$ and $\mathcal{X}_i$ is the domain of knob $k_i$.
Even for modest domains, $|\mathcal{X}|$ becomes combinatorially large.
In practice:
\begin{itemize}
    \item Many knobs are irrelevant to a target workload (e.g., USB settings for a database VM).
    \item Measurement budgets are limited: each trial configuration can require seconds to minutes of runtime.
    \item The objective function is noisy and non-stationary, resulting in a complicated exploration.
\end{itemize}

Thus, a naive autotuner that treats all knobs as candidates will be too slow, too costly, and often unstable.

\section{Research Hypothesis}\label{sec:problem-hypothesis}

This thesis is based on the following hypothesis:

\begin{quote}
\textbf{Hypothesis.} For a given application workload, only a small subset of Linux kernel knobs can actually influence performance; this subset can be approximated \emph{before} tuning begins by analyzing the kernel structure and the application--kernel interaction space.
\end{quote}

If the hypothesis holds, then an autotuner can focus on a reduced knob set $\mathcal{K}' \subset \mathcal{K}$ (with $|\mathcal{K}'| \ll |\mathcal{K}|$), improving sample efficiency and robustness.

\section{Problem Definition}\label{sec:problem-definition}

We now formalize the problem addressed in this thesis.

\subsection{System model}\label{sec:problem-model}

Instead of modelling a specific application-workload pair, the system model focuses on the operating system itself as a generalized platform for executing arbitrary server workloads. We assume a system characterized by:
\begin{itemize}
    \item \textbf{A target OS kernel} $L$ (e.g., a specific Linux kernel version) deployed in a standard server environment (bare-metal, virtual machine, or container).
    \item \textbf{Core kernel subsystems} $\mathcal{S}$ (such as the network stack, virtual memory manager, and CPU scheduler) that manage hardware resources, handle system calls, and dictate overall system throughput and latency.
    \item \textbf{A configuration space} $\mathcal{K}$, representing the full, high-dimensional set of runtime-tunable knobs exposed by $L$ (\texttt{sysctl}) that govern the operational policies of the subsystems in $\mathcal{S}$.
\end{itemize}

\begin{figure}[ht]
    \centering
    \begin{tikzpicture}[
        node distance=2.0cm, % Increased global spacing slightly
        box/.style={draw, rectangle, rounded corners, minimum width=2.5cm, minimum height=1cm, align=center, fill=white},
        arrow/.style={->, thick, >=stealth},
        % CHANGE 1: Increased 'inner sep' to 2.5em to push the label away from the boxes
        layer/.style={draw, dashed, inner sep=2.5em, label={[anchor=north west, font=\bfseries]north west:#1}}
    ]

    % --- USER SPACE ---
    \node[box, fill=blue!10] (agent) {\textbf{Tuning Agent}};
    \node[box, right=of agent, fill=gray!10] (cli) {\textbf{sysctl CLI}\\(procps)};
    \node[box, below=1.0cm of agent] (libc) {\textbf{C Library (glibc)}\\wrapper functions};
    
    % --- BOUNDARY ---
    % Adjusted line length to fit the wider spacing
    % \draw[thick, red] (-3, -2.8) -- (9, -2.8) node[pos=1, right] {\textbf{User/Kernel Boundary}};

    % --- KERNEL SPACE ---
    \node[box, below=2.5cm of libc, fill=red!10] (syscall) {\textbf{System Call Interface}\\(\texttt{sys\_sysctl})};
    \node[box, right=of syscall, fill=red!10] (proc) {\textbf{/proc/sys Handler}\\(procfs)};
    \node[box, below=1.2cm of syscall, fill=yellow!10] (data) {\textbf{Kernel Data Structures}\\(\texttt{tcp\_retries})};

    % --- ARROWS ---
    \draw[arrow] (agent) -- (libc) node[midway, right=3pt] {\small write()};
    \draw[arrow] (cli) |- (libc);
    \draw[arrow] (libc) -- (syscall) node[midway, right=3pt] {\small syscall / sysenter};
    
    % CHANGE 2: Added 'above=5pt' to push text up
    \draw[arrow] (syscall) -- (proc) node[midway, above=5pt] {\small dispatches};
    
    % CHANGE 2: Added 'right=5pt' to push text right
    \draw[arrow] (proc) -- (data) node[midway, right=5pt] {\small updates value};

    % --- BACKGROUND BOXES ---
    \begin{scope}[on background layer]
        \node[fit=(agent)(cli)(libc), layer=User Space, fill=blue!5] {};
        \node[fit=(syscall)(proc)(data), layer=Kernel Space, fill=red!5] {};
    \end{scope}

    \end{tikzpicture}
    \caption{Interaction between User-Space tuning agents and Kernel-Space data structures. The \texttt{sysctl} mechanism acts as a bridge, allowing user-space processes to modify kernel variables via the \texttt{/proc} virtual filesystem.}
    \label{fig:user_kernel_bridge}
\end{figure}

Although, downstream autotuners will eventually introduce a specific application $A$ and a workload distribution $W$ to measure real-world performance metrics, our foundational problem operates independently of $A$ and $W$.
The model relies solely on the internal architecture and source code of $L$ to deduce the structural relationship between the knobs in $\mathcal{K}$ and the performance-critical paths within $\mathcal{S}$.

\subsection{Optimization objective}\label{sec:problem-objective}

The classical autotuning objective is to find:
\begin{equation}
    \mathbf{x}^* \in \arg\max_{\mathbf{x} \in \mathcal{X}} f(\mathbf{x}; A, W, L, H)
\end{equation}

where the components of the equation are formally defined as follows:
\begin{itemize}
    \item $\mathbf{x}$: A specific configuration vector, representing a concrete assignment of values to the tunable kernel knobs.
    \item $\mathbf{x}^*$: the optimal configuration vector that yields the highest performance.
    \item $\mathcal{X}$: the high-dimensional search space of all possible valid configurations, bounded by the full set of available kernel knobs $\mathcal{K}$.
    \item $f(\cdot)$: the black-box objective function that measures system performance (e.g., throughput to be maximized, or the inverse of latency to be minimized). Evaluating this function requires executing a benchmark.
    \item $A$: the target application running on the system (e.g., a specific database, web server, or in-memory cache).
    \item $W$: the workload distribution driving the application (e.g., specific query mixes, read/write ratios, or request arrival patterns).
    \item $L$: the target OS kernel acting as the execution environment.
    \item $H$: the underlying hardware resources (e.g., CPU architecture, memory capacity, storage device) on which the system operates.
\end{itemize}

This optimization is typically subject to constraints such as safety limits, bounded exploration budgets (limiting the number of times $f$ can be evaluated), and operational policies.

This thesis does \emph{not} propose a new optimizer to evaluate or solve $f(\mathbf{x})$.
Instead, it addresses a prerequisite sub-problem that many optimizers implicitly require: reducing the dimensionality of the search space $\mathcal{X}$ before the tuning process begins. By mapping the entire space $\mathcal{X}$ to a much smaller, performance-critical subspace $\mathcal{X}'$ (derived from the reduced knob set $\mathcal{K}'$), we make the classical optimization objective computationally tractable for downstream autotuners.

\subsection{Knob selection (search-space reduction) problem}\label{sec:problem-knob-selection}

\textbf{Input:}
\begin{itemize}
    \item A Linux kernel version $L$,
    \item its source code and
    \item its full set of exposed runtime knobs $\mathcal{K}$.
\end{itemize}

\textbf{Output:} a universal, workload-agnostic subset of sysctl knobs $\mathcal{K}' \subseteq \mathcal{K}$ that contains parameters with a high potential to affect system performance across diverse workloads.
Tuning over $\mathcal{K}'$ instead of $\mathcal{K}$ preserves achievable performance gains while drastically reducing the dimensions exposed to the optimizer.

\begin{itemize}
    \item High \emph{recall}: effectively capture the critical knobs that influence core kernel performance paths.
    \item Good \emph{reduction factor (precision)}: aggressively filter out knobs that are strictly functional, debugging-related, or performance-neutral to mitigate the curse of dimensionality.
    \item \emph{Generality}: the resulting subset should serve as a strong starting baseline for any standard workload, delegating workload-specific fine-tuning to the downstream optimizer.
\end{itemize}

\section{Scope and Contributions}\label{sec:problem-scope}

To establish a clear boundary for this research, it is crucial to explicitly separate the core contributions of this thesis from the mechanisms used to evaluate them.

The strict focus of this work is the static identification and dimensionality reduction of kernel configuration parameters.
Specifically, the methodology targets runtime kernel knobs (such as \textit{sysctl} parameters) that can be dynamically altered without recompiling the operating system.
The primary contribution is a static analysis pipeline that leverages kernel code structure (mapping data flows and tracking variable usage within execution paths) to infer which specific parameters actively influence system behaviour.

Crucially, this methodology is completely decoupled from the runtime execution of any specific application.
It does not actively monitor, trace, or analyse the system call footprint of a live workload.
Instead, the pipeline simply accepts a predefined set of system calls as its input.
Given this input set, it statically traces the corresponding kernel execution paths to extract the relevant configuration parameters.
The system never needs to run an application, simulate traffic, or measure live telemetry to filter the search space.
By operating purely through offline source code analysis, it yields a highly targeted parameter subset based entirely on the intersection of the requested system calls and the kernel data flow.

\section{The Role of Autotuning (Non-Goal)}

It must be explicitly stated that designing, building, or improving an automated tuning optimization algorithm (such as a new Bayesian Optimization or reinforcement learning model) is strictly out of scope for this thesis.

In contrast, in this research an existing autotuning framework is utilized solely as a validation instrument.
By feeding the set of statically reduced parameters into an existing autotuner, we can empirically verify whether the identification pipeline successfully retained the most influential knobs while discarding the noise.
The autotuner serves only to demonstrate that our dimensionality reduction accelerates convergence and unlocks meaningful performance gains.

Other areas outside the scope of this research include compile-time kernel specializations (such as profile-guided optimizations or \textit{Kconfig} pruning) and the tuning of embedded or hard real-time systems, as our evaluation specifically targets representative cloud server environments.

\section{The Search Space Bottleneck}

The necessity of this static analysis approach becomes apparent when evaluating the current landscape of automated system tuning.
While the literature provides incredibly powerful optimizers, the prerequisite step of defining the search space remains a critical bottleneck.

The reliance on expert-curated list of parameters is labour-intensive, highly biased toward specific historical deployment scenarios, and rapidly becomes obsolete as kernel versions evolve.
Alternatively, workload-driven feature selection methods (such as Lasso~\cite{lasso} regression or Random Forests~\cite{random_forests}) are effective but operationally expensive.
They require a dynamic warm-up phase in which hundreds of benchmarking permutations must be executed against a live system simply to gather enough training data to rank the relevance of the knobs.
This massive cost defeats the purpose of a low-overhead tuning system.
Finally, attempts to automate selection via natural language processing (parsing kernel documentation or man pages) frequently fail because documentation is sparse and typically describes parameter functionality rather than actual performance impact.

This systemic gap motivates the proposed methodology: narrowing the configuration search space statically, without requiring expensive dynamic profiling runs or human domain expertise.

\section{Research Questions and Evaluation Criteria}

To structure the evaluation of this methodology, this thesis addresses the following core research questions:
\begin{enumerate}
    \item \textbf{Static identification}: can we effectively map user-facing \textit{sysctl} strings to their internal C variables and identify their performance relevance solely by analysing the static structure of the Linux kernel source code?
    \item \textbf{Dimensionality reduction}: can this static analysis successfully compress the vast configuration space, isolating a highly targeted subset of knobs?
    \item \textbf{Impact on tuning}: when an off-the-shelf autotuner is restricted to this statically derived subset, does the noise reduction lead to improved sample efficiency (requiring fewer benchmarking iterations to reach convergence) while achieving superior final throughput compared to tuning blindly across a generic knob set?
    \item \textbf{Evolutionary robustness}: is this extraction methodology robust across different kernel versions, automatically adapting to the introduction or deprecation of parameters without requiring manual rule rewriting?
\end{enumerate}

The success of the proposed system is evaluated against these questions using tangible metrics.
These include the percentage by which the search space is mathematically compressed and the application throughput achieved over the non-optimal kernel configuration.

\newpage
